\section{Background Modeling}
\label{sec:background-modeling}

The selected photon--jet invariant-mass spectrum is dominated by Standard
Model prompt-photon production in association with jets.  Smaller
contributions arise from multijet events in which a jet is misidentified as a
photon, with a further subleading contribution from electroweak
$W/Z+\mathrm{jets}$ production.  Since these processes produce a smoothly
falling spectrum and their
absolute rates are affected by sizeable theoretical and experimental
uncertainties, the background is determined directly from data rather than
from simulated samples.  This data-driven treatment also avoids relying on a
simulation-based prediction of the jet-to-photon misidentification rate.

Following previous searches for excited quarks in the photon+jet final state,
the background distribution is described by the four-parameter empirical
function
\begin{equation}
  \frac{\mathrm{d}N}{\mathrm{d}m_{\gamma j}}
  = P_{0}\,
    \frac{(1-x)^{P_{1}}}
         {x^{P_{2}+P_{3}\ln x}},
  \qquad
  x \equiv \frac{m_{\gamma j}}{\sqrt{s}},
  \qquad \sqrt{s}=13~\TeV ,
  \label{eq:background-function}
\end{equation}
where $P_{0}$ controls the normalization and $P_{1}$, $P_{2}$, and $P_{3}$
determine the shape.  This functional form is sufficiently flexible to model
the steeply falling spectrum while remaining smooth over the full search
region.

An independent background fit is performed for each tested $c^{*}$ mass
hypothesis.  The nominal peak position $x_{0}$ and the left and right widths,
$\sigma_{\mathrm{L}}$ and $\sigma_{\mathrm{R}}$, are taken from the
double-sided Crystal Ball signal model described in the preceding section,
with the signal-shape nuisance parameters fixed to their nominal values.  To
prevent a possible signal from biasing the
background estimate, the asymmetric window
\begin{equation}
  x_{0}-3\sigma_{\mathrm{L}}
  < m_{\gamma j} <
  x_{0}+3\sigma_{\mathrm{R}}
  \label{eq:background-excluded-window}
\end{equation}
is excluded.  A single fit to Eq.~\eqref{eq:background-function} is then made
simultaneously to the data in the lower and upper sidebands over
$700 < m_{\gamma j} < 3500~\GeV$.  Thus, the two sidebands constrain one common
set of parameters and one common covariance matrix; the fitted function is
used to interpolate smoothly across the excluded interval.  The window is
derived from the $f=1.0$ signal model at each mass and the same background
model is used for all three coupling hypotheses at that mass.

Figure~\ref{fig:background-fit} shows a representative fit for
$M_{c^{*}}=2.0~\TeV$.  The upper panel displays the data in the two sidebands
and the fitted function, while the vertical dashed lines indicate the excluded
signal window.  The lower panel shows the fractional residual,
$(N_{\text{data}}-N_{\text{fit}})/N_{\text{fit}}$.  No resonant structure is
introduced by the interpolation, and the residuals outside the masked region
are compatible with the smooth background description within the statistical
precision of the data.

\begin{figure}[tbp]
  \centering
  \includegraphics[width=0.82\textwidth]
    {/eos/user/h/hsiaoche/Data/Invariant_Mass_gJet_M2000.png}
  \caption{Fit of the four-parameter background function to the photon--jet
    invariant-mass distribution for the $M_{c^{*}}=2.0~\TeV$ hypothesis.  The
    data inside the interval bounded by the dashed lines are excluded from the
    fit.  The curve across this interval is the interpolation determined by a
    simultaneous fit to the lower and upper sidebands.  The lower panel shows
    the fractional difference between the data and the fitted prediction.}
  \label{fig:background-fit}
\end{figure}

For the statistical interpretation, a separate RooFit workspace is built for
each mass hypothesis.  The fit intervals are chosen to contain the relevant
signal line shape and sufficient sideband data: $500$--$1400~\GeV$ for the
$1.0~\TeV$ hypothesis, increasing progressively to
$2050$--$3500~\GeV$ for the $3.0~\TeV$ hypothesis.  Within each workspace,
$P_{1}$, $P_{2}$, and $P_{3}$ are unconstrained nuisance parameters and are
profiled in the likelihood fit.  The background normalization is also allowed
to float independently through a rate parameter.  Consequently, both the
normalization and the shape of the background are determined by the observed
data in the final fit, and the statistical uncertainty of the fitted
background is propagated to the confidence intervals.
